\setcounter{section}{1}


In der folgenden Aufgabe verwenden wir die Schreibweise

\mavergleichskette
{\vergleichskette
{D(r)
}
{ = }{ \left\{ (r,s,t) \in \R^3  \mid  r \neq 0        \right\}
}
{  }{
}
{    }{
}
{   }{
}
}
{}{}{.}

\inputaufgabe
{}
{





Bestimme zum
\definitionsverweis {Vektorbündel}{}{}

\mavergleichskettedisp
{\vergleichskette
{L
}
{ =} { \left\{ (r,s,t,u,v,w)  \mid  ru+sv+tw = 0 , \, (r,s,t) \neq (0,0,0)       \right\}
}
{ \subseteq} { \R^3 \setminus \{0,0,0\} \times \R^3
}
{ \longrightarrow} { \R^3 \setminus \{0,0,0\}
}
{ } {
}
}
{}{}{}
lineare Trivialisierungen oberhalb von
\mathl{D(r)}{,}
\mathl{D(s)}{} und
\mathl{D(t)}{,} also von $r,s,t$ abhängige Basen oberhalb von $D(r)$ usw. Bestimme die Basiswechselabbildungen auf

\mavergleichskette
{\vergleichskette
{D(rs)
}
{ = }{ D(r) \cap D(s)
}
{  }{
}
{    }{
}
{   }{
}
}
{}{}{.}





}
{} {}

\inputaufgabe
{}
{





Bestimme zum
\definitionsverweis {Vektorbündel}{}{}

\mavergleichskettedisp
{\vergleichskette
{L
}
{ =} { \left\{ (r,s,t,u,v,w)  \mid  ru+sv+tw = 0 , \, (r,s,t) \neq (0,0,0)       \right\}
}
{ \subseteq} { \R^3 \setminus \{0,0,0\} \times \R^3
}
{ \longrightarrow} { \R^3 \setminus \{0,0,0\}
}
{ } {
}
}
{}{}{}
diejenigen Parameter
\mathl{(r,s,t)}{,} für die der Vektor $\left(  3   , \,   7  , \,   4                 \right)$ zur Faser
\mathl{L_{(r,s,t)}}{} gehört.





}
{} {}

\inputaufgabe
{}
{





Zeige, dass in

Beispiel 1.2

auf
\mathl{D(r)}{} durch

\mavergleichskettedisp
{\vergleichskette
{u(r,s,t)
}
{ =} { { \frac{   t  }{ r } } \begin{pmatrix}  s  \\ -r\\  0   \end{pmatrix}  - { \frac{   s  }{ r } } \begin{pmatrix}  t  \\ 0 \\  -r   \end{pmatrix}
}
{ } {
}
{ } {
}
{ } {
}
}
{}{}{}
ein
\zusatzklammer {von den Parametern abhängiger} {} {}
Vektor im Lösungsraum definiert ist, der auf ganz $\R^3$ polynomial fortsetzbar ist, obwohl die Koeffizientenfunktionen
\mathkor {} {{ \frac{   t  }{ r } }} {und} {- { \frac{   s  }{ r } }} {}
nur auf $D(r)$ definiert sind und nicht fortsetzbar sind. Ist
\mathl{u(r,s,t)}{} überall Teil einer Basis?





}
{} {}

\inputaufgabe
{}
{





Bestimme in

Beispiel 1.3

die Parameter, für die der Lösungsraum
\mathl{L_{(a,b,c,d,e,f)}}{} ein-, zwei- oder dreidimensional ist. Sind diese Mengen offen oder abgeschlossen?





}
{} {}

\inputaufgabe
{}
{





Zeige, dass über einem beliebigen Körper $K$  zu
\definitionsverweis {linear unabhängigen}{}{}
Vektoren
\mathkor {} {u= \begin{pmatrix}  a  \\ b \\ c   \end{pmatrix}} {und} {v= \begin{pmatrix}  d  \\ e \\ f   \end{pmatrix}} {}
die Familie bestehend aus
\mathl{u,v}{} und dem
\definitionsverweis {Kreuzprodukt}{}{}

\mathl{\begin{pmatrix}  a  \\ b \\ c   \end{pmatrix}  \times \begin{pmatrix}  d  \\ e \\ f   \end{pmatrix}}{} keine
\definitionsverweis {Basis}{}{}
des $K^3$ bilden muss.





}
{} {}

\inputaufgabe
{}
{





Betrachte den
\definitionsverweis {topologischen Raum}{}{}

\mavergleichskettedisp
{\vergleichskette
{Y
}
{ \defeq} { \left\{ (s,t,u,v) \in \R^4  \mid  su+tv = 1        \right\}
}
{ } {
}
{ } {
}
{ } {
}
}
{}{}{}
mit der Projektion
\maabbeledisp {p} { Y } { \R^2 \setminus \{ (0,0) \} = X
} {(s,t,u,v)} { (s,t)
} {.}
\aufzaehlungvier{Zeige, dass jede Faser von $p$
\definitionsverweis {homöomorph}{}{}
zu einer reellen Geraden ist.
}{Zeige, dass durch

\mavergleichskettedisp
{\vergleichskette
{ \varphi(s,t)
}
{ =} {(s,t,u(s,t),v(s,t))
}
{ =} { \left( s,t, { \frac{   s  }{ s^2+t^2 } } , { \frac{   t  }{ s^2+t^2 } }  \right)
}
{ } {
}
{ } {
}
}
{}{}{}
eine stetige Abbildung
\maabb {\varphi} { X} { Y
} {}
mit

\mavergleichskettedisp
{\vergleichskette
{ p \circ \varphi
}
{ =} {
\operatorname{Id}_{  X  }
}
{ } {
}
{ } {
}
{ } {
}
}
{}{}{}
gegeben ist.
}{Definiere einen
\definitionsverweis {Homöomorphismus}{}{}
zwischen
\mathkor {} {Y} {und} {X \times \R} {.}
}{Zeige, dass es keine polynomiale Abbildung
\maabb {\psi} { X } { Y
} {}
mit

\mavergleichskettedisp
{\vergleichskette
{ p \circ \psi
}
{ =} {
\operatorname{Id}_{  X  }
}
{ } {
}
{ } {
}
{ } {
}
}
{}{}{}
gibt.
}





}
{} {}

\inputaufgabe
{}
{





Zeige, dass ein
\definitionsverweis {reelles Vektorbündel}{}{}
über einem Punkt
\zusatzklammer {also einem einpunktigen topologischen Raum} {} {}
das gleiche ist wie ein reeller endlichdimensionaler Vektorraum.





}
{} {}

\inputaufgabe
{}
{





Es sei
\maabb {p} { V } { X
} {}
ein
\definitionsverweis {reelles Vektorbündel}{}{}
über einem
\definitionsverweis {topologischen Raum}{}{}
$X$. Zeige, dass $V$ genau dann ein
\definitionsverweis {Hausdorffraum}{}{}
ist, wenn $X$ ein Hausdorffraum ist.





}
{} {}

\inputaufgabe
{}
{





Es sei $X$ ein
\definitionsverweis {topologischer Raum}{}{.}
Zeige, dass man die
\definitionsverweis {Identität}{}{}
\maabb {\operatorname{Id}_{  X  }} { X } { X
} {}
als ein
\definitionsverweis {reelles Vektorbündel}{}{}
vom
\definitionsverweis {Rang}{}{}
$0$ auffassen kann.





}
{} {}

\inputaufgabe
{}
{





Es sei $X$ ein
\definitionsverweis {topologischer Raum}{}{.}
Zeige, dass ein
\definitionsverweis {Homomorphismus}{}{}
der
\zusatzklammer {trivialen} {} {}
Vektorbündel
\maabbdisp {\varphi} {X \times \R^n } { X \times \R^m
} {}
das gleiche ist wie eine
$m \times n$-\definitionsverweis {Matrix}{}{,}
der Einträge stetige Funktionen von $X$ nach $\R$ sind.





}
{} {}

\inputaufgabe
{}
{





Es sei

\mavergleichskette
{\vergleichskette
{U
}
{ \subseteq }{ \R^n
}
{  }{
}
{    }{
}
{   }{
}
}
{}{}{}
\definitionsverweis {offen}{}{}
und
\maabb {f} { U } { \R^m
} {}
eine
\definitionsverweis {stetig differenzierbare Abbildung}{}{.}
Zeige, dass durch das
\definitionsverweis {totale Differential}{}{}
in der Form
\maabbeledisp {} { U \times \R^n } { U \times \R^m
} { (x,v)} { (x, \left(D f \right)_{ x} (v) )
} {,}
ein Homomorphismus des Vektorbündels
\mathl{U \times \R^n}{} in das Vektorbündel
\mathl{U \times \R^m}{}  gegeben ist.





}
{} {}

\inputaufgabe
{}
{





Zeige, dass es eine
\definitionsverweis {Homöomorphie}{}{}
des
\definitionsverweis {Tangentialbündels}{}{}

\mathl{T_{S^1 }}{} der
$1$-\definitionsverweis {Sphäre}{}{}
$S^1$ mit dem
\definitionsverweis {Produkt}{}{}
$S^1 \times \R$ gibt.





}
{} {}
Was hat die vorstehende Aufgabe mit

Beispiel 1.1

zu tun?

\inputaufgabe
{}
{





Man gebe ein Beispiel einer differenzierbaren Kurve
\maabbdisp {\gamma} { [0,1[} { S^1
} {}
derart, dass der
\definitionsverweis {Grenzwert}{}{}

\mathl{\operatorname{lim}_{   t  \rightarrow  1   } \,  \gamma(t)}{} existiert, dass aber der Grenzwert
\mathl{\operatorname{lim}_{   t  \rightarrow  1   } \, ( \gamma(t), T_t (\gamma ) (1))}{} in
\mathl{TS^1}{} nicht existiert.





}
{} {}

\inputaufgabe
{}
{





Zeige, dass die Abbildung
\maabbeledisp {} {TS^1} { \R^2
} {( (a,b),t (-b,a))} { (a,b) + t (-b,a)
} {,}
für jeden Punkt
\mathl{(x,y) \in \R^2}{} außerhalb der Einheitskreisscheibe zwei Urbildpunkte, auf dem Einheitskreis einen Urbildpunkt und innerhalb der offenen Einheitskreisscheibe keinen Urbildpunkt besitzt. Man interpretiere dies geometrisch.





}
{} {}

\inputaufgabe
{}
{





Man gebe ein Beispiel einer
\definitionsverweis {injektiven}{}{}
\definitionsverweis {differenzierbaren Abbildung}{}{}
\maabbdisp {\varphi} { M } { N
} {}
zwischen zwei
\definitionsverweis {differenzierbaren Mannigfaltigkeiten}{}{}
\mathkor {} {M} {und} {N} {}
derart, dass die zugehörige
\definitionsverweis {Tangentialabbildung}{}{}
\maabbdisp {T(\varphi)} {TM} {TN
} {}
nicht injektiv ist.





}
{} {}

\inputaufgabe
{}
{





Man gebe ein Beispiel einer
\definitionsverweis {surjektiven}{}{}
\definitionsverweis {differenzierbaren Abbildung}{}{}
\maabbdisp {\varphi} { M } { N
} {}
zwischen zwei
\definitionsverweis {differenzierbaren Mannigfaltigkeiten}{}{}
\mathkor {} {M} {und} {N} {}
derart, dass die zugehörige
\definitionsverweis {Tangentialabbildung}{}{}
\maabbdisp {T(\varphi)} {TM} {TN
} {}
nicht surjektiv ist.





}
{} {}

\inputaufgabe
{}
{





Seien
\mathkor {} {M} {und} {N} {}
\definitionsverweis {differenzierbare Mannigfaltigkeiten}{}{} und es sei
\maabb {\varphi} { M } { N
} {} eine
\definitionsverweis {differenzierbare Abbildung}{}{.} Zeige, dass die zugehörige
\definitionsverweis {Tangentialabbildung}{}{}
\maabbdisp {T(\varphi)} {TM} {TN
} {}
\definitionsverweis {stetig}{}{}
ist.





}
{} {}
